\section{First joins of previously constructed tables}\label{sec:joins}

The direct comparisons in Section~\ref{sec:direct} cover adjacent
simultaneous-birth pairs. We now bound all remaining first-sharing pairs,
including nonadjacent pairs born at the same height. The principal difficulty is that the old
weights depend on earlier matchings. We retain their actual prefixes in a
backward calculation. A covariance term is then charged to the two
positions that tested the corresponding deletions, using the capacities
of the earlier matching matrices.

\begin{proposition}[Remaining first-join comparisons]
\label{join:remaining-first-joins}
Consider a finite block and the retained centres, masks, and table
operations of Sections~\ref{sec:centres}--\ref{sec:partitions}, with the
stated parameter hierarchy. At each first sharing of a component, omit
the simultaneous-birth comparisons already treated in
Section~\ref{sec:direct}. For the remaining comparisons, the following
estimates hold uniformly over the permitted prior table history.
\begin{enumerate}
\item The sum of broad pair counts in $\lambda\otimes\lambda$, at
distance at most $C r_{\max}$ and with width weight $r_{\min}$, is $O(1)$.
Here $C$ is any fixed positive constant.
\item The same-width, same-remaining-mask counts in common new cells,
weighted by the single factor $r_i\pi_I(a)$ and the product of the two
raw weights, have total
\[
 O(W)+O_{k_0}(W^{1+c})
\]
for some fixed $c>0$.
\end{enumerate}
The constants are independent of the slot cap $N_0$ and of the finite
block. Thus, together with the direct comparisons, these estimates
establish Proposition~\ref{tab:first-join-input}.
\end{proposition}

We prove the proposition in stages. All point counts below are per unit
length, with one first-point lift chosen and the second point ranging
over the relevant near lifts. We use the centre notation
\[
 G=(L,M),\qquad U=L/G,\qquad T=M/G,\qquad
 \Theta_{ij}=G r_{\min},\qquad K=[L,M]r_{\max}.
\]
When a bound is summed over masks, its outer-mask coefficients are kept
separate: they are independent probability factors for a broad count,
and one common probability factor for a clipping count.

\subsection{The two kinds of first join}

\begin{lemma}[Classification of the remaining pairs]
\label{join:classification}
Suppose centres $i,j$ first share a component at height $Y$, and their
comparison is not a simultaneous-birth comparison treated in
Section~\ref{sec:direct}. Either
\begin{equation}\label{join:distant}
 \log(UT)>\beta Y/2,
\end{equation}
or, after interchanging the centres, $i$ is old, $j$ is a medium-type
newborn, and
\begin{equation}\label{join:mixed}
 Y=A_0 B_j.
\end{equation}
In this case $B_i/B_j$ is a sufficiently small fixed fraction,
$\log(UT)=O(\beta A_0 B_j)$, and the sum of the
reciprocals of all prime divisors of both centres below bin $Y$ is at most
$s_{\rm low}B_j$, with the arbitrarily small constant specified in the
parameter hierarchy. The same alternatives apply after replacing an old
endpoint by a centre in one of its earlier components.
\end{lemma}

\begin{proof}
If both centres are newborn but their pair was not treated by the
direct comparisons, then $\log(UT)>\beta Y$: those comparisons cover
the adjacent newborn pairs. Such a pair satisfies
\eqref{join:distant}, even if a path of other newborn centres puts
the two centres in one component.

If both centres are old, they belonged to different components at the
preceding height $Y/2$. The previous adjacency threshold therefore gives
\eqref{join:distant}. This also defines the distant case when one or both
centres are newborn.

Suppose \eqref{join:distant} fails. The two centres are then directly
adjacent at activation. A high-type newborn has the same last dense bin
as every neighbour at that height, since an exponent discrepancy at a
prime in that bin already exceeds the edge log-cost. Such neighbours
start together, so they give a simultaneous-birth comparison. The
remaining newborn is consequently medium type. Initially its start is
at most $2D_0 B_j$ up to the fixed onset convention, and hence
$\log(UT)=O(\beta D_0B_j)$. The centre exponents agree at every prime
whose bin has height at least $\ell B_j$.

Centres in the same band would have the same start. A centre in a larger
retained band cannot start earlier. Thus the old centre belongs to a
smaller band, separated from $B_j$ by the prescribed large factor. If a
common dense bin existed at height at least $\ell B_j$, its height would
exceed $D_0 B_i$, contradicting the retained medium type of $i$. There is
therefore no such dense bin. In particular the newborn starts at
$A_0B_j$, rather than at twice a later dense height. The bins below
$\ell B_j$ contribute $O(\ell B_j)$ to the reciprocal-prime sum; all
later bins below $Y$ are sparse and contribute $O(\kappa A_0 B_j)$.
The choices of $\ell$ and $\kappa A_0$ give the assertion.

An earlier component of an old centre is contained in its old child at
the current join. Replacing that centre inside the earlier component
keeps it old and keeps the two endpoints in different old children, or
keeps it opposite the same newborn. The preceding argument applies to
the replaced pair as well.
\end{proof}

Call the alternatives \emph{distant} and \emph{mixed}, respectively.
For a comparison at height $Y$, fix its processed masks $d,e$ and write
\[
 v'=L/d,\qquad w'=M/e.
\]
All prime factors of $d,e$ are below the present bin. Write $a,b$ for
the outer masks, with either their independent laws or their single
common law. Here $I$ is the component's current common remaining part:
the broad coefficient is $\pi_I(a)\pi_I(b)$, whereas clipping requires
$a=b$ and uses just $\pi_I(a)$. Set
\[
 f_0=(Td,Ue),\qquad
 H=B_i+B_j+\log(2UT),\qquad P_0=\rho H,
\]
where the fixed constant $\rho>0$ is sufficiently small.

At relative distance at most $S r_{\max}$, the determinant integer lies
in an interval of length $O(SK/f_0)$. Its cycle has width-weighted size
\begin{equation}\label{join:bare-cycle}
 r_{\min}(v',w')=\Theta_{ij}\frac{f_0}{de}.
\end{equation}
The bare main term is consequently $O(S k_L k_M/(de))$. In a broad
comparison $S$ is a fixed constant. In the clipping comparison the
preliminary partition tests permit
\begin{equation}\label{join:narrow-distance}
 S\ll \exp\bigl(-E_g\max(B_i,B_j)\bigr),\qquad r_i=r_j.
\end{equation}
We drop actual common-cell membership after using this distance bound.
This removes the history-dependent cell predicate from all the
subsequent counts.

For a projected denominator $v'$, put
\[
 R_{\rm done}(v')=
 \prod_{\substack{l\mid v'\\l\ \mathrm{processed}}}R_l.
\]
The identity
\begin{equation}\label{join:processed-currency}
 \frac{R_{\rm done}(v')}{d}=\pi_{\rm done}(d)
\end{equation}
uses the deficit law on the prime powers already processed at this
join. We will return the main terms to this currency, preserving the
outer coefficients separately.

\subsection{A backward expansion of the actual weights}

For an old side, its \emph{adaptive range} consists of the primes from
its start up to, but not including, the present bin. The adaptive range
of a newborn is empty. Above $P_0$, the reciprocal sum of mismatch
primes dividing $UTde$ is bounded uniformly. For $UT$ this follows from
its log-size; for $d,e$ it follows from the good-mask reciprocal purge,
since $d_0\ll\rho$. In particular, for every fixed $C$,
\begin{equation}\label{join:mismatch-restoration}
 \prod_{\substack{p>P_0\\p\mid UTde}}R_p^{-C}\ll 1.
\end{equation}
The constant can depend on the fixed cutoffs, but not on the centres or
masks in the block.

A processed prime $p>P_0$ is called a \emph{bad omission} for the fixed
pair if a side adaptive at $p$ omits it and either the two centre
exponents at $p$ differ or the other side omits $p$ as well.

\begin{lemma}[Backward decomposition]
\label{join:backward-decomposition}
Read in decreasing order every prime divisor of either centre above
$P_0$ whose size lies in the union of the two adaptive ranges. For the
fixed current grids, masks, and distance
predicate, the weighted pair count is bounded by the sum of the
following nonnegative terms:
\begin{enumerate}
\item a continuing product term, stopped at the low fields or at the
first bad omission;
\item one separately stopped joint-loss term at each common full factor
where both deleted parts are nonconstant and at least one is an
adaptive linked deletion.
\end{enumerate}
On the continuing term, already read primes supply, up to a uniform
constant, the product of the corresponding ordinary $R_p$ factors on
the two rows. Each adaptive joint-loss term has the actual raw link
masses as coefficients. No condition inherited from a larger-prime
alignment remains on the continuing term.
\end{lemma}

\begin{proof}
Keep the pair domain on its current grid throughout this expansion.
The only position restriction is the distance predicate; row and mask
restrictions remain allowed. On a side adaptive at a prime $p$, denote
the actual ancestor weights just before and just after its update by
$w_x^-$ and $w_x^+$. Between successive primes read backwards, discard
any intervening clipping, preliminary tests, or further multipliers
from the upper bound. Monotonicity then bounds the current weight by
$w_x^+$, and after processing the update the continuing term retains
$w_x^-$. If the current row omitted $p$, its locations already lie on
the successful ancestor subgrid; no unrestricted success probability
is inserted at this point.

When this procedure leaves a side's adaptive range, its actual birth
weight is an upper bound for the remaining prefix. Replace that weight
by the product upper version of Proposition~\ref{tab:virtual-marginals}.
In each fixed static anchor choice, keep one active birth exclusion at
each prime, chosen for that vertex independently of its position or
partner; if no exclusion is active, keep the actual scalar. Drop the
other filters. Later-processed masks use primes disjoint from the birth
digits and hence multiply the current numerator by units at all these
static primes. The retained exclusions are therefore fixed residues
in the current numerator. Average the static anchor choices only after
the corresponding upper calculation, independently for the two
vertices.

If both adaptive ranges are empty, make these static replacements
immediately. There are no bad-omission or joint-loss terms in
that case. The continuing term goes directly to the remaining-field
count, which retains all high static factors above its low cutoff.
This includes nonadjacent newborn pairs in the distant case.

Stop the continuing term at its first bad omission. At any other
mismatch, whether a centre-exponent discrepancy or an actual omission,
drop the field at $p$. Equation~\eqref{join:mismatch-restoration} pays
for restoring its ordinary factor throughout this branch.

It remains to describe a common full factor $p^\nu$. Simultaneous
translation by $1/p^\nu$ preserves the two current grids and their
distance. On adaptive sides it preserves the actual pre-$p$ prefixes
and transports the links equivariantly. On a static side the smaller
retained factors are invariant. Write the adaptive multiplier as
$1-\mathcal D_p$, where $\mathcal D_p$ is the linked-success deletion
fraction plus $1/p$ times the unpaired fraction. Its orbit mean is
$1/p$. A static field is either a fixed forbidden colour or the
constant $R_p$.

More explicitly, at an adaptive full-factor point $x$, the centered
deleted raw mass is
\begin{equation}\label{join:centered-link-loss}
 \sum_{\xi}\mathfrak m_{x,\xi}
 \left(\mathbf 1_{\{\xi\ \mathrm{succeeds}\}}-\frac1p\right),
\end{equation}
where $\mathfrak m_{x,\xi}$ is the actual raw mass linked from $x$ to
the plus entry $\xi$. These masses are transported unchanged on the
joint orbit. Indeed, writing
$M_x=\sum_\xi\mathfrak m_{x,\xi}\le w_x^-$ and
$I_\xi=\mathbf1_{\{\xi\ \mathrm{succeeds}\}}$, the deleted raw mass is
\[
 \frac{w_x^--M_x}{p}+\sum_\xi\mathfrak m_{x,\xi}I_\xi.
\]
Consequently
\[
 w_x^+=R_pw_x^-
       -\sum_\xi\mathfrak m_{x,\xi}(I_\xi-1/p).
\]
For two adaptive full-factor sides, orbit averaging gives the exact
identity
\begin{equation}\label{join:raw-covariance}
 \mathbb E(w_x^+w_y^+)
 =R_p^2w_x^-w_y^-
 +\sum_{\xi,\eta}\mathfrak m_{x,\xi}\mathfrak m_{y,\eta}
       \left(\mathbb E(I_\xi I_\eta)-\frac1{p^2}\right).
\end{equation}
The centered part of a static forbidden field has the
same form with one indicator; a scalar field has no centered part.
Expanding the product of the two updates, the orbit average is at most
$R_p^2$ times the two pre-prefixes plus the average joint term of their
nonconstant success-losses. Indeed the covariance expansion of
\eqref{join:centered-link-loss} contains the joint-success kernel minus
the nonnegative product-of-means kernel.

Continue only the product-main term. Stop each joint-loss term at
this prime, retaining the actual link masses and, for a static
indicator, its smaller static factors. At least one side is adaptive
at a prime in the union being read, so a nonzero separately stopped
term has an adaptive linked side. Summing over locations turns the
joint orbit averages into the corresponding counts of actual
successes. Such a stopped term is not subsequently expanded into
further adaptive covariance terms.

This construction never assumes independent match plans. It retains
the actual prefix whenever a link decision is needed, and each
nonconstant coefficient is precisely $\mathfrak m_{x,\xi}$, not a
fractional coefficient multiplied by an enlarged prefix. All the
ordinary factors from larger read primes are now constants. Their
mismatch restoration is uniformly bounded by
\eqref{join:mismatch-restoration}, proving the stated form of every
branch.
\end{proof}

The expansion has three kinds of terminal term: the continuing product
reaches the low fields, it stops at its first bad omission, or a separate
joint-loss term stops where it splits off. The first two keep the
original centre pair and require a periodic count of the remaining
fields. The third requires an additional calculation before it can be
charged to the positions that tested the corresponding deletions.

\subsection{Counting terms without adaptive links}

\begin{lemma}[Remaining-field count]\label{join:endpoint-count}
The retained fields through $P_0$ on the two projected grids have
periods dividing those denominators, of sizes at most
$\exp(O(P_0))$, and their separate product means are the corresponding
ordinary factors. Remaining high static fields, if present, increase
the main-term constant by a bounded factor and the roundoff by at most
$\exp(C\Sigma)$, where
\[
 \Sigma=\sum_{\substack{p\mid LM\\ \log\log p<Y}}\frac1p.
\]
The period loss can be taken at most $\exp(\epsilon H)$ by choosing
$\rho$ sufficiently small after $\epsilon>0$ is fixed.
\end{lemma}

\begin{proof}
For adaptive low fields, the modulus and separate product means follow
from Corollary~\ref{tab:intermediate-fields} and
Equation~\eqref{eq:core-loss}. If only static factors remain, their prime
residue periods give the same bound directly.
Apply Lemma~\ref{ar:cycle} to the product of the two retained low-field
weights. Their separate means give the independent residue mean in that
lemma, and the product of their periods is at most $\exp(O(P_0))$.

Suppose some high static factors remain. Then their low companion
fields are static and have coprime moduli. Drop high mismatch factors
at the uniformly bounded restoration cost. At a high common full
factor, the two forbidden residues either coincide on the determinant
cycle or do not. Their cycle density, relative to the two ordinary
factors, is at most
\begin{equation}\label{join:static-coincidence}
 \prod_p\left(1+\frac Cp
       \mathbf 1_{\{n\equiv n_p\pmod p\}}\right),
\end{equation}
where $n$ is the determinant integer and $n_p$ is the possible
coincidence residue. With only one static factor at $p$ there is no
extra factor. The digit CRT separates these high factors from the low
fields.

Expand \eqref{join:static-coincidence}. A term imposes one congruence
on $n$ modulo a squarefree product $Q'$. Since $Q'$ is coprime to the
low periods, the progression form of Lemma~\ref{ar:cycle} gives the interval length
divided by $Q'$, plus $O(\exp(O(P_0)))$ roundoff, both multiplied by
the ordinary mean of the low fields. The main extra product has local
increments $O(1/p^2)$ and is bounded. The roundoff extra product is at
most $\exp(C\Sigma)$.

There is no unhandled combination of high static fields and adaptive
low fields. If either side is adaptive at or below $P_0$ on the
continuing term, the union of adaptive ranges covers every higher
prime below the join; those factors have already been read. Replace
the remaining actual prefix by its product upper version through
$P_0$ on the current grid. If a bad omission has stopped the branch
and lower adaptive fields are dropped, only static residue fields are
needed. Finally $P_0=\rho H$ and a sufficiently small $\rho$ give the
claimed period loss.
\end{proof}

Consider first the continuing term when it reaches the low fields.
Lemma~\ref{join:endpoint-count} restores
$R_{\rm done}$ on both sides in the main term. By
\eqref{join:processed-currency}, summing the masks with their assignment
restrictions gives $O(m_Lm_M)$. In the clipping count the small factor
$S$ in \eqref{join:narrow-distance} absorbs the possible inverse
assignment density arising from the common outer mask.
To see this explicitly, for side $i$ let
\[
 A_i(a)=\sum_{\substack{d:\,da\text{ assigned and good}}}
               \pi_{{\rm done},i}(d).
\]
Then $0\le A_i\le1$ and $\mathbb E_a A_i\le\delta_{L_i}$.
Independent outer laws give at most $\delta_L\delta_M$. For one
outer law,
\[
 \mathbb E_a(A_i(a)A_j(a))\le\min(\delta_L,\delta_M),\qquad
 \frac{S\mathbb E_a(A_iA_j)}{\delta_L\delta_M}
 \le\frac{S}{\max(\delta_L,\delta_M)}\ll1,
\]
by $\delta_{L_i}\ge e^{-.64B_i}$ and the choice of $E_g$.

Now suppose the term stops at a bad omission $p$, in bin $Y_p$. Drop
adaptive cuts at $p$ and below, while keeping the upper ordinary
factors already obtained. Keep the static upper fields through this
point and apply Lemma~\ref{join:endpoint-count}. For this bound use a
broad distance even in the clipping count. Since the adaptive range
starts in sparse bins, restoring the dropped adaptive factors costs
at most $\exp(C\kappa Y_p)$.

Designate one adaptive omitted side, say $L$. Drop its assignment
restriction. Summing its processed mask gives a factor $1/p$ while
the other side's assignment remains in force. This assertion is
uniform in the fixed outer mask: the unrestricted processed law has
$\Pr(p\mid d)=1/p$. Thus for the common outer law the remaining sum
is at most $\mathbb E_a A_j(a)/p\le\delta_M/p$, and the same bound
holds with independent outer laws. If $p\mid UT$, then
$p>P_0$ and
\[
 \sum_{\substack{p\mid UT\\p>P_0}}\frac{\log p}{p}
 \le \frac{\log(UT)}{P_0}\ll_\rho1.
\]
Since $\log\log p\ge A_0B_i$ on the designated adaptive side and
$\exp(C\kappa Y_p)\le(\log p)^{C\kappa}$, this implies
\begin{equation}\label{join:bad-prime-saving}
 \sum_{\substack{p\mid UT,\ p>P_0\\
                  \log\log p\ge A_0B_i}}
 \frac{\exp(C\kappa Y_p)}p
 \ll \exp(-cA_0B_i).
\end{equation}
If instead $p$ divides the partner mask $e$, the same bound holds at
each fixed good partner mask. Indeed $p>P_0>d_0B_j$, and its good-mask
bound controls $\sum_{p\mid e}(\log p)^{.3}/p$ in this range. Choose
$C\kappa<.3$ to obtain the exponential saving from the same lower
cutoff. This saving pays for $\delta_L^{-1}$ in
\eqref{eq:centre-lower}, also for a coupled outer law. The resulting
main terms again sum to $O(m_Lm_M)$: before this last restoration
their bound is
$k_Lk_M\delta_M e^{-cA_0B_i}
 =m_Lm_M\delta_L^{-1}e^{-cA_0B_i}$.

For either kind of terminal term just described, the roundoff per fixed mask choice,
with its outer coefficients still outside, is at most
\begin{equation}\label{eq:no-link-roundoff}
 C\Theta_{ij}\frac{f_0}{de}
 \exp(\epsilon H+C\Sigma)
 \bigl(1+\omega(d)+\omega(e)\bigr).
\end{equation}
Here density gains have been dropped. The omega factor allows every
possible bad stopping prime; there is only one continuing branch.

\subsection{Charging joint-loss terms to successful plus entries}

At a separately stopped term at $p>P_0$, move each adaptive linked
full-factor portion fractionally along its link to the plus entry
whose success removed it. This transfers an upper-bound counting
predicate; it does not move any retained table location. Denote the
image centres by $L',M'$, leaving an unmapped static centre unchanged.
On a mapped first side, let $d'$ be the below-$p$ processed deficit of
the image plus row, whose pre-success denominator is $L'/d'$; use $e'$
analogously on a mapped second side. These image deficits need not
equal the original minus deficits $d_{<p},e_{<p}$.
Here $d_{<p},e_{<p}$ are the parts of the original masks $d,e$
supported below $p$.

Each image has the same width and belongs to the same component as its
origin at the $p$-step. Each origin and its image have the same prime powers from $p$
upward, and the matching preserves the entire mask after $p$. Cell
span allows us to enlarge the distance predicate to $O(r_{\max})$ on
the image pair. On a static partner, retain its current-grid forbidden
residue at $p$ and its smaller static factors instead.

At this stage it is only the distance predicate that has moved. The
original full-factor locations still satisfy the divisibility
conditions imposed by their later projections. We first average those
original-grid restrictions, then sum the origins using matching
capacities, and finally count the image grids and restore their
first-mass factors. The following exact averaging supplies the first
of these steps.

\begin{lemma}[Averaging later omissions]\label{join:later-omissions}
For each mapped side, write
\[
 d=d_{<p}b_i^+,
\]
where $b_i^+$ is the processed portion of its actual mask above $p$
and below the join bin. In the linked comparison, expansion to the
full minus grid at the $p$-step, followed by averaging its later
projection congruences, costs exactly $1/b_i^+$. These factors multiply
over the mapped sides.
\end{lemma}

\begin{proof}
The full ancestor on this side has denominator $L/d_{<p}$, phase
$b_i^+a\gamma$, and numerator $b_i^+z$, where $z$ is the current
numerator. The post-success plus image has the same phase. Expand back to
this minus grid while retaining the condition that its numerator is
zero modulo $b_i^+$. This is exactly the original current projection,
not a condition on the linked image.

Let $l\mid b_i^+$. All these later omissions are adaptive. Since the
branch did not stop at a larger bad omission, the two original centres
have the same exponent $\nu_l$ and the partner has no omission at $l$.
Translate both sides and all participating links by multiples of
$1/l^{\nu_l}$. Both expanded grids are closed under these
translations, including the current grid of an unmapped static
partner. The translations preserve the pre-$p$ prefix data and links,
success at $p$, the retained factors through $p$ on a static side,
and the broadened distance between the image locations. They also
preserve the congruences at other later-omission primes.

More precisely, put $V_0=L/d_{<p}$, so $\nu_l(V_0)=\nu_l$.
The designated numerator increment is $V_0/l^{\nu_l}$, a unit
modulo $l$ and a multiple of every other later-omission prime and
of $p$. On the pre-success plus grid the denominator is $L'/d'$,
which also has its full $p$-power; its numerator increment under
the same translation is therefore a multiple of $p$, preserving
the success residue. The partner is full at $l$, so it has the same
grid-closure property. Thus all the other congruences and $p$-success remain
unchanged, whereas the designated numerator runs uniformly modulo $l$. Its
divisibility condition therefore averages to $1/l$. Good adaptive
masks have deficit at most one at these primes, and omissions above
$p$ on two mapped sides are disjoint. Successive primewise averaging
gives the claimed product. Equivalently, the translations permute
tuples consisting of minus points and their links, with any static
forbidden entry. All upper ordinary factors are already constants.
No survival condition from a later prime or partition remains in
these tuples, so the image need not survive any later projection.
\end{proof}

\begin{lemma}[Capacity and uniqueness of the charge]
\label{join:capacity}
After the averaging of Lemma~\ref{join:later-omissions}, summing all
originating full-factor labels charges at most the weight of each
successful image plus entry. With two mapped endpoints the bound is
the product of these capacities. The charged image pair first shares
a component at the same join as the original pair, and there is no
extra sum over origin masks or comparison epochs. This holds for
both the product-law broad count and the one-law clipping count.
\end{lemma}

\begin{proof}
Fix an image plus label at its $p$-step, its width, and its post-$p$
mask. Its incoming link masses sum to at most its own raw weight,
and hence to at most one. The fixed post-$p$ mask, split at the
current join height, determines both
$b_i^+$ and the current outer mask $a$. Every originating centre in
this $p$-component has the same prime powers from $p$ upward.
Consequently its carried upper ordinary factors, the factor
$1/b_i^+$, and its current outer-law coefficient depend only on these
preserved data, not on the individual preimage.

Below-$p$ origin masks are already labels in this one link matrix;
they are included in its column sum. Moreover each involved minus
ancestor already has an original assigned-row label, including its
entire future mask. It therefore has only one possible current row
and mask labelling, with its later numerator fixed by the projection
condition before that condition was averaged. Thus no additional processed-mask sum occurs
after the capacity bound. If both sides are mapped, widening the
domain to a predicate on the image pair alone permits the two column
bounds to be multiplied. If a side is static, keep its same upper
static choice throughout this summation and then average its anchors.
Other origin restrictions may be discarded.

In particular, for two fixed successful image labels $\xi,\eta$
with the same outer mask $a$, let $c_\xi,c_\eta$ denote their
preserved ordinary and later-omission factors. Nonnegativity and
the column capacities give
\begin{equation}\label{join:one-law-capacity}
 \begin{split}
 &\sum_{x,y}\pi_I(a)c_\xi c_\eta
           \mathfrak m_{x,\xi}\mathfrak m_{y,\eta}\\
 &\qquad\le
 \pi_I(a)c_\xi c_\eta
       \left(\sum_x\mathfrak m_{x,\xi}\right)
       \left(\sum_y\mathfrak m_{y,\eta}\right)
 \le \pi_I(a)c_\xi c_\eta w_\xi w_\eta.
 \end{split}
\end{equation}
The original pairs can be a restricted subset of the displayed
Cartesian product. There is one outer coefficient on both sides
of this inequality, not its square. Independent outer laws give
the same calculation with $\pi_I(a)\pi_I(b)$ instead.

The earlier $p$-component of each mapped centre is contained in the
appropriate old child of the current join. Hence an image pair could
not have shared a component earlier: that would already have joined
the two old children. In a mixed comparison the newborn is unchanged
and has no earlier table. Thus each image pair is charged at its own
first-sharing epoch. The choice of mapped sides has only finitely
many possibilities. Finally the outer coefficient is preserved,
whether it is $\pi_I(a)\pi_I(b)$ or the single $\pi_I(a)$.
\end{proof}

Bound the mapped plus raw weights by one, retaining the smaller static
factors on any unmapped side. The joint-loss term is then bounded by a
count on two grids with denominators
\[
 \frac{L'}{pd'},\qquad \frac{M'}{pe'}.
\]
On a mapped side $d'$ is the below-$p$ processed mask of the new plus
row. On a static side put $e'=e$; its forbidden residue at $p$ turns
the current grid into one of denominator $M/(pe)$ with a modified
phase. This transformation multiplies the smaller static residues by
a $p$-unit and translates them, so their counts and means remain
unchanged. There is at most one static side in this branch.

Primed centre parameters now refer to the image pair. The
path log-costs to the image centres are
$\ll\exp(h_*Y_p/10)$ by \eqref{eq:component-path}, whereas a mapped
centre has band at most $Y_p/A_0$. The original inequality $p>\rho H$
therefore implies
\[
 p>\rho'H'
\]
for some fixed $\rho'>0$. Use a smaller smoothing cutoff against
$H'$ in Lemma~\ref{join:endpoint-count}. Any high mismatches of the
image pair restore at bounded relative cost by its good masks, after
taking $d_0$ small enough. Retained static fields need no adaptive
averaging here.

The image determinant cycle has width-weighted size
\[
 \Theta'\frac{(T'd',U'e')}{pd'e'},
\]
and its determinant interval has length
\[
 O\left(\frac{[L',M']r_{\max}}
                  {p(T'd',U'e')}\right).
\]
Its main term is therefore
\begin{equation}\label{join:linked-main}
 \frac{k_{L'}k_{M'}}{p^2d'e'},
\end{equation}
multiplied by the carried upper ordinary factors, the factors
$1/b_i^+$ on the mapped sides, and the ordinary factors supplied by
the remaining static tests.

For a mapped side the first-mass quantity to restore is
\[
 k_{L'}\pi_{\rm done}(d'pb_i^+).
\]
Its geometric factor $k_{L'}/(pd'b_i^+)$ and its ordinary factors
above $p$ are exactly present in \eqref{join:linked-main}. Explicitly,
if $D'$ is the processed prime-power part of $L'$ at the join, then
\[
 \pi_{D'}(d'pb_i^+)
 =\frac{R(D'/(d'pb_i^+))}{pd'b_i^+}.
\]
Only its ordinary factors through $p$ remain to be restored, at cost
at most $\exp(CY_p)$. Drop this side's assignment restriction and
sum the deficit law: the required deficit one at $p$ has unrestricted
probability at most $1/p$. Thus the geometric $1/p$ becomes the
forced-omission probability; it does not disappear in this sum. On a
static partner keep the actual assignment restriction on $e'$; its
forbidden $p$-residue supplies an additional $O(1/p)$. The same
calculation works with a coupled outer mask, because mapped sides
have dropped their assignment restrictions while the static side
retains its own. At a fixed common outer mask $a$, the mapped
processed sum is at most $1/p$ and the static side contributes at
most $A_j(a)/p$ after its field means are restored. Averaging $a$
gives at most $\delta_{M'}/p^2$. Two mapped sides instead give
at most $1/p^2$ before restoring their assignment densities.

For each mapped centre, $Y_p\ge A_0B'$. Summing the resulting
$1/p^2$ over these primes more than pays both $\exp(CY_p)$ and the
inverse assignment densities from \eqref{eq:centre-lower}. Thus all
such main terms are charged by $O(m_{L'}m_{M'})$ to the image pair.
This conclusion uses the capacities and unique comparison epoch of
Lemma~\ref{join:capacity} before summing centre pairs.

The roundoff, per image masks and stopping prime, is at most
\begin{equation}\label{eq:linked-roundoff}
 C\Theta'\frac{(T'd',U'e')}{pd'e'}
 \prod_{\mathrm{mapped}\ i}(b_i^+)^{-1}
 \exp(\epsilon H'+C\Sigma').
\end{equation}
We have dropped ordinary-factor gains, but kept the outer coefficients.
Unlike the main-term restoration, this bound has no
$\exp(CY_p)$ loss.

\subsection{Summing roundoff and completing the comparisons}

All main terms are now bounded by products of first masses. It remains
to sum the two roundoff bounds \eqref{eq:no-link-roundoff} and
\eqref{eq:linked-roundoff}. The charged pair is the original pair in
the first bound and the image pair in the second. Write its parameters
without primes in either formula. At common centre exponents the
local unrestricted mask sum for the gcd ratio is $1+O(1/p)$; at an
exponent discrepancy it is bounded by a fixed divisor power times
$1+O(1/p)$, as in Section~\ref{sec:direct}. The same calculation
permits the factor $2^{\omega(d)+\omega(e)}$, which bounds the
omega factor in \eqref{eq:no-link-roundoff}. Additional upper masks
sum reciprocally, and in \eqref{eq:linked-roundoff} the stopping
primes are summed with their factor $1/p$. These sums give
\begin{equation}\label{eq:summed-roundoff}
 C\Theta_{ij}\tau(UT)^C
       \exp(\epsilon H+C\Sigma).
\end{equation}
The outer laws, including the one-law version, are probability laws
throughout. The coefficients of $\Sigma$ in this estimate are
absolute.

\paragraph{Distant pairs.}
Here $\Sigma\ll Y$ and \eqref{join:distant} holds. Take $\beta$
sufficiently large and then $\epsilon$ sufficiently small. Substituting
\eqref{eq:centre-kernel} and \eqref{eq:centre-lower} into
\eqref{eq:summed-roundoff}, the factor $\exp(CY)$ is paid by a small
part of the available negative power of $UT$. With mass exponent
$u=.75$, a strict margin remains for the tensor estimate, as does a
fixed residual exponential decay in $Y$. The factor $\tau(UT)^C$
costs an arbitrarily small further power. The residual decay pays the
shell multiplicities and the height sum, while a positive power of
the width ratio remains for summing exact dyadic widths. These
roundoffs consequently total $O_{k_0}(W^{1+c})$ for some $c>0$.

\paragraph{Mixed pairs.}
Let $j$ be the newborn larger-band centre. By
Lemma~\ref{join:classification}, $\log(UT)=O(\beta A_0B_j)$ and
$\Sigma$ is arbitrarily small relative to $B_j$. If
\[
 \log(r_i/r_j)>.65 B_j^*,
\]
the width-gain consequence following \eqref{eq:exterior-gain} applies
to \eqref{eq:summed-roundoff}. With sufficiently small harmonic and
period exponents, it leaves an exponential saving in $B_j$ and a
positive residual width power, again summable with a mass power
beyond one.

Otherwise the exterior alternative in \eqref{eq:exterior-gain}
holds for every assigned row of $j$. Keep its processed mask $c_j$
fixed and set $W'=M/c_j$. Its good-mask condition gives
$\tau(c_j)\le e^{s'B_j}$, since the entire processed segment has the
small reciprocal-prime sum. Divide \eqref{eq:no-link-roundoff} by
this row's first-mass quantity
$r_jW'R_{\rm done}(W')$, leaving its outer law in place. Summing
only the old side's processed mask gives
\begin{equation}\label{join:mixed-relative-error}
 C\exp(\epsilon H+C\Sigma)\frac{r_{\min}}{r_j}
 \left(\frac{W'}{(W',L)}\right)^{-1}
 \tau(U)^C\tau(c_j)^C.
\end{equation}
For completeness, before this sum the geometric ratio is
$\frac{r_{\min}}{r_j}(W',L/d)/W'$. Its sum is bounded exactly as in
the conditional target-grid calculation of
\eqref{eq:target-grid}; the extra omega factor is absorbed by fixed
divisor moments. The inverse ordinary density costs at most
$\exp(C\Sigma)$.
Writing $E_{\rm ext}=W'/(W',L)$, the primewise sum just used is
\[
 \sum_d\frac{(W',L/d)}{W'}
 \ll E_{\rm ext}^{-1}\tau(U)^C\tau(c_j)^C\exp(C\Sigma).
\]
At a prime, the free deficit allowance before the denominator of
this ratio grows is at most $\nu_p(U)+\nu_p(c_j)$; beyond that
allowance the terms decrease geometrically. This proves the bound
also with the fixed divisor moments for the omega factor.

In a mixed joint-loss term the newborn side is necessarily static;
its centre and assigned row are unchanged. Its
fictitious restricted denominator is $W'/p$, and the other counted
denominator also contains the factor $1/p$. Relative to the original
newborn first mass this gives the same bound
\eqref{join:mixed-relative-error}, with an additional $1/p$ whose
sum is absorbed in $\exp(C\Sigma)$. The exterior ratio is unchanged
by dividing this common full factor from both denominators: at the
stopping prime $p$ the two image centres have the same full exponent
and $p$ divides neither $d'$ nor $c_j$, so
\[
 \left(\frac{L'}{pd'},\frac{W'}p\right)
   =\frac1p(L'/d',W'),\qquad
 \frac{W'/p}{(W'/p,L'/p)}=\frac{W'}{(W',L')}.
\]
The classification of Lemma~\ref{join:classification} still applies
to the image pair, so its old centre has the required smaller band.
Dropping origin restrictions did not change the newborn assigned
row or its divisibility by $W'$.

An assigned true denominator $v$ of $j$ divides $W'$. Therefore
\[
 \frac{W'}{(W',L)}\ge \frac{v}{(v,L)},
\]
and \eqref{eq:exterior-gain} makes the product of the width factor
and the inverse exterior ratio in
\eqref{join:mixed-relative-error} exponentially small in $B_j$.
Choose first $s'$ small enough for $\tau(c_j)^C$, then the small
harmonic allowance, and finally the period and divisor-function
powers sufficiently tiny relative to $\beta A_0$. There remains an
exponential saving more than sufficient to sum all smaller-band
centres $i$ by their centre count. The resulting contribution is
$O(W)$, and can be made a small fixed multiple of $W$ where needed.

All terminal main terms were bounded by the product of the original or
charged image masses, so their total is $O(W^2)$. The distant and
width-separated roundoffs are $O_{k_0}(W^{1+c})$, and the oriented
mixed roundoffs are $O(W)$. These conclusions apply separately to
the broad and narrow counts with their respective mask laws. Since
$W\le1$, they prove Proposition~\ref{join:remaining-first-joins}.
In particular none of the constants used for clipping depends on
$N_0$.

\subsection{Closure of the finite construction}

We finish by explaining the order in which the proved estimates are
used. This also records the parameter margins needed to make the
finite construction uniform over arbitrarily late blocks.

\begin{proof}[Proof of Theorem~\ref{thm:main}]
The elementary initial reduction deals with widths for which
$vr(v)$ does not tend to zero. In the remaining case fix a putative
positive-measure set avoided by a whole tail, dyadically reduce the
positive widths, and use the finite late blocks of mass between
$w$ and $2w$ described in Section~\ref{sec:preliminaries}.

The row and centre selections and the birth prescriptions are fixed
from the block before the adaptive history starts. The additional
birth holes in Section~\ref{sec:direct} depend only on those assigned
lists. They therefore do not require any estimate for an already
running table. Direct comparisons establish the adjacent simultaneous-birth
part of Proposition~\ref{tab:first-join-input}.

Proceed through the finitely many component starts, component joins,
and sparse-prime steps in chronological order. At a component change,
the histories from earlier steps and their deterministic danger
trains are already defined. Section~\ref{sec:partitions} constructs
the new padded partitions from these histories, proving
Proposition~\ref{tab:partition-input} with negligible first-mass
loss. The preliminary tests give the cell spans used in the present
section. Proposition~\ref{join:remaining-first-joins} then supplies
the remaining broad and narrow comparisons for the actual old
prefixes, before the new clipping operation. The narrow comparison
bound permits that clipping. At each following sparse-prime update,
Section~\ref{sec:sparse} applies to the existing component and
partitions; it supplies Proposition~\ref{tab:sparse-input} and the
new danger trains for later component changes. Its nonexceptional
tight pairs coarsen the accounting classes as prescribed. Thus no
stage needs an estimate depending on a future partition or a future
adaptive matching. Induction completes the finite construction.

We recall the quantitative margins. In sparse bins the harmonic
losses that compete with a height saving are at most
$\exp(C(\sigma+\kappa)Y)$, with absolute $C$, whereas the off-gap
window supplies decay with exponent $(1-\sigma_g)Y$. The choices of
$\sigma,\kappa$ leave a strict margin. The switching budgets in
the sparse and direct comparisons are taken sufficiently large for
their crude unneutralized costs; sparse errors may carry
$\tau(UT)^C$ separately. At foreign-child partition tests and
distant first joins, the coefficients of $C\Sigma$ or $CY$ are
absolute, so the initial large choice of $\beta$ pays them. At
mixed comparisons the divisor loss needs only absolute smallness
of $Cs'$. After these choices the period and exterior powers can be
taken as small as required relative to $\beta A_0$.

Take $k_0$ small enough to absorb all bounded-height effects and to
make the exponentially decaying selection and history losses small
relative to the selected mass and the fixed avoided-set measure.
The colour, type, and hard-mean constants have already been fixed.
The later choice of the large fixed threshold $\lambda_0$ makes
the birth-hole losses small. Its holes use squarefree low moduli,
so this choice does not enlarge the prime powers required for core
equivariance. Finally choose $w$ sufficiently small for the terms
$O_{k_0}(W^{1+c})$ and choose the fixed cap $N_0$ sufficiently large
for the narrow comparison losses. All these choices are independent
of the location of the finite block in the tail.

We have now proved all three inputs to
Proposition~\ref{tab:conditional-reduction}. The energy bound of
Proposition~\ref{tab:energy-bound}, the virtual marginals of
Proposition~\ref{tab:virtual-marginals}, and their equidistribution
give the contradiction to the positive avoided set exactly as in
that reduction. Consequently, for the fixed shift and width
sequence, almost every point belongs to infinitely many of the
prescribed interval trains. This proves the theorem.
\end{proof}
